% from ega0-22-fr.tex
\oldpage[0\textsubscript{IV-1}]{207}

discrètes (19.3.5, (iv)), $B_{\mathfrak p'}$, muni de la topologie
$\mathfrak p'$-préadique, est encore une $A_{\mathfrak q}$-algèbre
formellement lisse lorsque l'on prend sur $A_{\mathfrak q}$ la topologie
discrète ou la topologie $\mathfrak q$-préadique (19.3.8). Pour montrer
l'équivalence de b) et d) lorsque $B$ est noethérien, appliquons (22.6.2) à
l'anneau $A_{\mathfrak q}$ discret (ou $\mathfrak q$-préadique) et à l'anneau
$B_{\mathfrak p'}$, muni de la topologie $\mathfrak p'$-préadique~; comme
$C_{\mathfrak p}$ est préadmissible et que $\mathfrak pC_{\mathfrak p}$ est
un idéal de définition pour sa topologie, on peut invoquer l'équivalence de
c) et a) dans (22.6.2).
\end{proof}

\begin{corollary}[22.6.5]
\label{0.22.6.5-fr}
Sous les hypothèses générales de (22.6.4), l'ensemble des
$\mathfrak p\in\operatorname{Spec}(C)$ tels que $C_{\mathfrak p}$ soit une $A$-algèbre
formellement lisse (ou une $A_{\mathfrak q}$-algèbre formellement lisse, en
désignant par $\mathfrak q$ l'image réciproque de $\mathfrak p$ dans $A$)
pour les topologies discrètes, est ouvert dans $\operatorname{Spec}(C)$.
\end{corollary}

Cela résulte de la forme c) de (22.6.4). On notera que lorsque $B$ est
noethérien, on peut remplacer les topologies discrètes par les topologies
préadiques.

\begin{corollary}[22.6.6]
\label{0.22.6.6-fr}
Sous les hypothèses générales de (22.6.4), les conditions suivantes sont
équivalentes~:
\begin{enumerate}
  \item[a)] $C$ est une $A$-algèbre formellement lisse (pour les topologies
  discrètes).
  \item[b)] Pour tout $\mathfrak p\in\operatorname{Spec}(C)$ (ou seulement pour tout idéal
  maximal $\mathfrak p$ de $C$), $C_{\mathfrak p}$ est une $A$-algèbre
  formellement lisse (ou une $A_{\mathfrak q}$-algèbre formellement lisse, en
  désignant par $\mathfrak q$ l'image réciproque de $\mathfrak p$ dans $A$)
  pour les topologies discrètes.
\end{enumerate}
En outre, lorsque $B$ est noethérien, on peut remplacer dans b) les topologies
discrètes par les topologies préadiques sur $A_{\mathfrak q}$ et
$C_{\mathfrak p}$.
\end{corollary}

\begin{proof}
La dernière assertion résulte de (22.6.4). En vertu de (22.6.3), la condition
a) équivaut au fait que l'homomorphisme (22.6.1.1) est inversible à gauche~;
l'équivalence de a) et b) résulte donc de (19.1.14), compte tenu du fait que
$\Omega^1_{B/A}$ est un $B$-module projectif et $\mathfrak J/\mathfrak J^2$
un $C$-module de type fini.
\end{proof}

\begin{proposition}[22.6.7]
\label{0.22.6.7-fr}
Soient $k_0$ un corps, $k$ une extension séparable de $k_0$, $C$ une
$k$-algèbre de type fini.
\begin{enumerate}
  \item[\rm(i)] Soit $\mathfrak p$ un idéal premier de $C$. Les conditions
  suivantes sont équivalentes~:
  \begin{enumerate}
    \item[a)] $C_{\mathfrak p}$ est une $k_0$-algèbre formellement lisse pour
    la topologie discrète.
    \item[b)] $C_{\mathfrak p}$ est une $k_0$-algèbre formellement lisse pour
    la topologie $\mathfrak p$-préadique.
    \item[c)] Il existe $f\in C-\mathfrak p$ tel que $C_f$ soit une
    $k_0$-algèbre formellement lisse pour la topologie discrète.
    \item[d)] $C_{\mathfrak p}$ est un anneau géométriquement régulier
    (19.6.5) sur $k_0$.
  \end{enumerate}
  En outre l'ensemble des $\mathfrak p\in\operatorname{Spec}(C)$ ayant l'une quelconque de
  ces propriétés est ouvert dans $\operatorname{Spec}(C)$.
  \item[\rm(ii)] Les conditions suivantes sont équivalentes~:
  \begin{enumerate}
    \item[a)] $C$ est une $k_0$-algèbre formellement lisse pour la topologie
    discrète.
    \item[b)] Tout $\mathfrak p\in\operatorname{Spec}(C)$ vérifie les conditions
    équivalentes de (i).
    \item[c)] Tout idéal maximal $\mathfrak m$ de $C$ vérifie les conditions
    équivalentes de (i).
    \item[d)] Pour toute extension $k'$ de $k_0$, $C\otimes_{k_0}k'$ est un
    anneau régulier (17.3.6)~; on dit encore que $C$ est un anneau
    géométriquement régulier sur $k_0$.
  \end{enumerate}

  